\section{Goals}\label{sec:goals}

\S\ref{sec:evidence-logic} introduced the two axes on which the evidence logic judges an agent's
work --- \emph{goals} (properties of what was established) and \emph{policies} (constraints on how it
was established) --- and gave the goal half in outline: a goal is a formula over \trace{} whose atoms
range over artifacts, lineage, and freshness, resolved against the evidence rather than self-reported.
This section deepens that half. We give the goal's construction, the three orthogonal verdicts it
carries, the typed acceptance criterion and the fuller property language it is a fragment of, the
lineage-based resolution rule and its non-monotonic honesty overlay, and an honest account of what is
authorable today versus what remains to be built.

\subsection{Constructing a goal}\label{sec:goals-construction}

A goal is the \emph{definition of done}, declared once, up front, and thereafter a passive lens. It is
declared through a single operation --- \code{declare\_goal} --- that records five things:
the \code{intent} (the user's words, restated), the \code{scope} (the files and symbols the goal
touches), optional \code{intent\_clauses} (the intent split into individually checkable propositions),
the \code{acceptance} items (the criteria that constitute ``done''), and optional inline
\code{policies} (the governance bar; see \S\ref{sec:policies}). The declared goal is embedded directly
into \trace{} as a first-class field (\code{goals}, canonicalizing to the empty list when no intent is
declared), so that ``definition of done'' and ``progress'' are computed by the same shared resolver
the rest of the substrate uses --- not by a private, per-tool notion. The declared goal's shape is
shown in Figure~\ref{fig:goal-goal-type}.

\begin{figure}[H]
\centering
\small
\[
\id{goal} \;::=\; \Big(\;
\begin{array}{@{}l@{\;:\;}l@{}}
\id{goal\_id}        & \id{String},\\
\id{intent}          & \id{String} \quad\text{\footnotesize\itshape (the change and why, in plain language)},\\
\id{intent\_author}  & \opt{\id{Authorship}} \quad\text{\footnotesize\itshape (who stated the intent, usually Human)},\\
\id{intent\_clauses} & \many{\id{String}} \quad\text{\footnotesize\itshape (intent split into checkable clauses)},\\
\id{scope}           & \many{\id{String}} \quad\text{\footnotesize\itshape (files / symbols the goal touches)},\\
\id{acceptance}      & \many{\id{AcceptanceItem}} \quad\text{\footnotesize\itshape (the end node: what ``done'' means)},\\
\id{policies}        & \opt{\id{GoalPolicies}} \quad\text{\footnotesize\itshape (the governance bar, \S\ref{sec:policies})},\\
\id{criteria\_review} & \opt{\id{CriteriaReview}} \quad\text{\footnotesize\itshape (the ``is it right?'' sign-off)},\\
\id{status}          & \opt{\id{GoalStatus}},\\
\id{meta\_action\_id} & \opt{\id{String}} \quad\text{\footnotesize\itshape (the meta-action pursuing this goal)}
\end{array}
\;\Big)
\]
\caption{The \code{goal} type recorded by \code{declare\_goal} and embedded in \trace{} as a first-class field, capturing intent, scope, acceptance criteria, policies, and the criteria-review sign-off.}
\label{fig:goal-goal-type}
\end{figure}

Two construction properties matter downstream, both consequences of the goal being decoupled from the
reasoning that satisfies it. First, re-declaring is \emph{idempotent}: re-stating a goal carries over
any already-resolved criterion rather than wiping a partly-proved one. Second, because the goal is a
formula over \trace{}'s evidence and \emph{not} the evidence itself, it can be attached after the fact
and is therefore \emph{portable} --- a team's definition of done is a small, versioned, human-authored
artifact that, once imported, re-resolves against \emph{your} trace with no evidence transferred. What
must line up is only the referenced components: a glob scope travels cleanly, while an exact symbol
reference needs that symbol present (or the durable identity of \S\ref{sec:identity} to rebind it
across a rename).

\subsection{Three orthogonal axes}\label{sec:goals-axes}

Where \S\ref{sec:evidence-logic} named \code{met}, \code{governed}, and \code{certified} as the
verdicts, here we make them precise. Trustworthy done is the conjunction of all three, and each is
independent and separately reported.

\begin{center}
\footnotesize
\begin{tabular}{@{}llll@{}}
\toprule
Axis & Question & Mechanism & Principal \\
\midrule
\code{met}       & Does the required \emph{artifact} exist?        & criterion $\leftrightarrow$ artifact by type + lineage & the doer's evidence \\
\code{governed}  & Was it produced compliantly?                    & the goal's policies over its cone      & org / context (a pack) \\
\code{certified} & Is the definition of done itself \emph{right}?  & criteria review by a non-doer          & a party $\neq$ the doer \\
\bottomrule
\end{tabular}
\end{center}

The axes are orthogonal \emph{by construction}, and the separation is deliberate. \code{met} asks only
whether the evidence \emph{exists}; it says nothing about quality. A refuted \code{VerificationResult}
still exists, so a criterion demanding a verification artifact for a component is met by it --- whether
that result had to be \emph{proved} rather than left refuted is a \code{governed} question, decided by
a policy over lineage, not a criterion field. This is why a goal can be \emph{met but not governed}:
the required artifact is present, but it was too weak (or produced non-compliantly) for the goal's bar
--- a state that today reads silently as ``done'' and here becomes first-class and surfaced. The
\code{certified} axis is orthogonal in a different direction: it asks not whether the formula holds but
whether it is the \emph{right} formula, and its normative rule is that the party who \emph{meets} a
goal must not be the sole party who \emph{decides its definition of done} --- so certification is a
sign-off by a principal other than the doer. Because how the work was done cannot change what was
established, \code{met} $\perp$ \code{governed} is a theorem we machine-check
(Table~\ref{tab:proofs}, \emph{governed $\perp$ met}).

\subsection{The acceptance criterion and the property language}\label{sec:goals-criterion}

An acceptance criterion, in the shipping encoding, is a typed record pairing a \emph{code component}
with the \emph{evidence artifact} that serves as its warrant (Figure~\ref{fig:goal-acceptance-criterion}).
That is all it encodes --- it does not say whether the artifact is good enough.

\begin{figure}[H]
\centering
\small
\[
\id{evidence\_req} \;::=\; \Big(\;
\begin{array}{@{}l@{\;:\;}l@{}}
\id{artifact} & \id{String} \quad\text{\footnotesize\itshape (the artifact \textsc{type} required in the component's lineage:}\\
              & \text{\footnotesize\itshape \quad $\id{VerificationResult} \mid \id{Decomp} \mid \id{Tests} \mid \id{Diff} \mid \id{Documentation} \mid \ldots$)}
\end{array}
\;\Big)
\]
\[
\id{acceptance\_criterion} \;::=\; \Big(\;
\begin{array}{@{}l@{\;:\;}l@{}}
\id{id}        & \id{String},\\
\id{statement} & \id{String} \quad\text{\footnotesize\itshape (human ``what must be true'')},\\
\id{component} & \big(\id{function\_} : \id{String},\; \id{file} : \opt{\id{String}}\big),\\
\id{evidence}  & \id{evidence\_req},\\
\id{required}  & \id{Bool},\\
\id{covers}    & \many{\id{String}} \quad\text{\footnotesize\itshape (intent clauses it addresses)}
\end{array}
\;\Big)
\]
\caption{The shipping acceptance-criterion encoding: a typed record binding a code component to the evidence-artifact type required as its warrant, together with the intent clauses it covers.}
\label{fig:goal-acceptance-criterion}
\end{figure}

This fixed struct is a \emph{concrete encoding of a fragment} of a fuller \textbf{property language} ---
the evidence logic proper --- and every legacy form desugars into it, so no existing trace breaks. The
property language expresses a goal's ``done'' as a conjunction of \emph{obligations}, each a predicate
over \trace{} tagged with the verdict axis it feeds. The atoms are exactly the resolver predicates the
implementation already computes, listed in Figure~\ref{fig:goal-atoms}.

\begin{figure}[H]
\centering
\small
\[
\begin{array}{@{}l@{\qquad}l@{}}
\id{has}(c,T)      & \text{an artifact of type $T$ roots in component $c$ (latest wins)}\\
\id{proved}(c)     & \id{has}(c,\id{VerificationResult}) \land \id{status} \in \{\id{proved},\id{sat}\}\\
\id{refuted}(c)    & \id{has}(c,\id{VerificationResult}) \land \id{status} = \id{refuted}\\
\id{conforms}(S)   & \text{a $\id{ConformanceResult}$, complete, $\id{diverged}=0$, over sources $S$}\\
\id{decomposed}(c) & \text{a decomposition roots in $c$}\\
\id{tested}(c)     & \text{a region-complete / coverage-rooted suite roots in $c$}\\
\id{no\_open\_gap}(c) & \text{no open residual touches $c$}\\
\id{fresh}(e)      & \text{evidence $e$ is closure-fresh (\S\ref{sec:state}, I4)}\\
\id{uncontested}(e)& \text{$e$ is not contested by an open defeater}
\end{array}
\]
\caption{The atoms of the property language --- the resolver predicates the implementation already computes over \trace{}, from artifact existence and verdict checks to freshness and contestation.}
\label{fig:goal-atoms}
\end{figure}

\noindent\textbf{Selectors} turn coverage into a quantifier rather than one row per symbol:
\code{component(f)}, \code{symbol(s)}, \code{module(p)}, \code{glob("payments/**")}, \code{tag(l)},
\code{scope}, and \code{high\_stakes} (the configured high-stakes glob set). An obligation is then
\code{[role]\ formula}, where \code{formula} combines atoms with \code{and}, \code{or}, \code{not},
\code{=>} and the quantifiers \code{forall $x$ in $sel$} / \code{exists $x$ in $sel$}, and
\code{role} \(\in\) \code{\{[met], [governed]\}} (default \code{[met]}). The role is not
``deliverable versus standard'': a \code{[met]} obligation is an \emph{outcome} property over the
resulting artifacts (\code{proved($c$)}, \code{conforms($S$)}, \code{no\_open\_gap($c$)}), while a
\code{[governed]} obligation is a \emph{process} constraint over how those artifacts were produced
(\code{formalize before verify}, \code{eventually reproved}). The connectives and quantifiers are
shared; the \emph{sort ranged over} --- the artifact sort versus the action sort --- \emph{is} the
goal-versus-policy distinction of \S\ref{sec:evidence-logic}. Freshness is a \emph{modality}, not a
truth condition: an obligation that holds but whose witness fails \code{fresh} is
\emph{true-but-\code{at\_risk}}, not false --- the honesty layer, made first-class in the language.

The desugaring makes one subtlety explicit. A typed \code{\{component, evidence\}} criterion is
\code{has($c$, T)} --- \emph{existence}, so a refuted result still satisfies it --- whereas the legacy
\code{property} kind is \code{proved($c$)}, which checks the verdict and so resolves a refuted result to
\code{blocked}. In the language one simply picks the atom that means what is intended, as the
desugarings in Figure~\ref{fig:goal-desugaring} show.

\begin{figure}[H]
\centering
\small
\[
\begin{array}{@{}l@{\;\rightsquigarrow\;}l@{}}
\big(\id{component}{:}\{\id{function\_}{:}f\},\; \id{evidence}{:}\{\id{artifact}{:}T\}\big) & [\id{met}]\; \id{has}(f, T)\\[2pt]
\id{kind}{:}\,\id{property}\ \{\id{symbol}\}      & [\id{met}]\; \id{proved}(\id{symbol})\\[2pt]
\id{kind}{:}\,\id{obligation}\ \{\id{policy\_id}\} & [\id{governed}]\; \id{holds}(\id{policy\_id})\\[2pt]
\id{kind}{:}\,\id{gap}\ \{\id{residual\_id}\}      & [\id{met}]\; \id{no\_open\_gap}(\id{target\text{-}of}\ \id{residual\_id})\\[2pt]
\id{kind}{:}\,\id{change}\ \{\id{symbol}\}         & [\id{met}]\; \id{has}(\id{symbol},\; \id{Diff} \lor \id{IMLModel})
\end{array}
\]
\caption{How each legacy criterion kind desugars into a role-tagged atom of the property language, so no existing trace breaks and one picks the atom that means what is intended.}
\label{fig:goal-desugaring}
\end{figure}

\subsection{Resolution by lineage}\label{sec:goals-resolution}

Resolution \emph{evaluates} the formula against the trace; nothing is authored as done. A criterion
\code{has($c$, T)} (and its verdict-checking cousins) resolves \code{met} when \trace{} holds an
artifact of type \code{T} whose lineage \textbf{roots in} the component $c$ --- read from structured
provenance (\code{target\_symbol}, \code{derived\_from}, the model's symbol list), never from a
description string. This is the fix for the ``goal never ticks'' defect: binding by a substring of the
goal's prose silently misses when the proved symbol's name and the description's wording diverge, so the
rooting is computed over lineage instead. When several artifacts match, the latest by producer action is
the evidence pointer. Concretely the resolver binds as shown in Figure~\ref{fig:goal-match}.

\begin{figure}[H]
\centering
\small
\[
\begin{array}{@{}l@{}}
\id{match}(o) \;=\; \operatorname*{argmax}_{a \in \mathcal{C}}\; \id{producer\_action\_id}(a), \quad\text{where}\\[4pt]
\mathcal{C} \;=\; \{\, a \;\mid\; \id{artifact}(a) \land \id{type}(a) = T \land \id{roots\_in\_component}(a, c) \land \lnot\,\id{contested\_by\_open\_defeater}(a) \,\}
\end{array}
\]
\caption{The resolver's binding rule: among artifacts of the required type whose lineage roots in the component and that are not contested by an open defeater, the latest by producer action becomes the evidence pointer.}
\label{fig:goal-match}
\end{figure}

\noindent then writes onto the \emph{enriched copy} (never back onto the authored goal) a \code{status}
(\code{done} on a match, \code{blocked} on a matched-but-refuted-or-contested result, \code{todo} on no
match), the matched artifact's id as \code{evidence\_ref}, and the \code{at\_risk} flag with a reason
when the witness is stale or detached.

\textbf{Component precision.} \code{roots\_in\_component} is \emph{target-precise}: an artifact that
declares its own \code{target\_symbol} (a verification goal, a decomposition, a targeted diff) is about
\emph{that} component, not every symbol the shared model happens to formalize. Only when nothing in the
lineage names a target does resolution fall back to the model's symbol list. Without this, a
decomposition of one function would spuriously appear to root in every symbol of the enclosing model.
Because the match is by component-plus-type over lineage rather than by a stored pointer, one artifact
can satisfy the same obligation across \emph{several} goals --- goals are lenses over one shared evidence
stream, not containers owning separate traces --- and a rename is followed
(\code{roots\_in\_component} accepts the source name \emph{or} the IML symbol; \S\ref{sec:identity}).

\textbf{Progress and the atom-to-helper mapping.} Progress is the fraction of required items resolved:
\(\text{progress} = |\{\text{resolved}\}| / |\{\text{total}\}| \in [0,1]\) (with an in-progress item
counting as a half), and a goal is \code{met} exactly when \emph{all} required items resolve. Each atom
delegates to a shipping resolver helper --- \code{proved}/\code{has} to the typed-lineage resolver over
\code{roots\_in\_component}, \code{no\_open\_gap} to the residual scan, \code{fresh}/\code{uncontested}
to the freshness verdict and the open-defeater check. These facts are machine-checked
(Table~\ref{tab:proofs}, \emph{goal resolution}): \emph{progress} \(\in [0,1]\),
\emph{met $\Leftrightarrow$ all-done}, and \emph{at\_risk never demotes}.

\textbf{The \code{at\_risk} honesty overlay (non-monotonic).} Because the trace is a non-monotonic
context (\S\ref{sec:context}, I4), a met obligation can drift. When the evidence a met item rests on
goes stale --- its target's dependency-closure content changed --- the item \emph{stays met} but is
flagged \code{at\_risk}, with a reason and a pointer to the derived stale-evidence residual. It never
silently demotes to unmet, and it never silently passes as clean: the two are distinguished, and a
re-established, fresh result clears the flag. A resolved obligation captures the closure checksum its
witness was verified against, so an edit trips exactly the same freshness authority that governs the
rest of the substrate. Figure~\ref{fig:goal-at-risk} traces a met obligation through a source edit and
a superseding re-proof.

\begin{figure}[H]
\centering
\small
\noindent\textit{Example (met but \id{at\_risk}).} The obligation $[\id{met}]\;\id{proved}(\id{refund})\;[\id{fresh}]$ is
resolved at three checkpoints:
\[
\begin{array}{@{}l@{\;\rightarrow\;}l@{}}
\text{A: \id{refund} proved fresh} & \id{done},\; \id{evidence\_ref} = \id{VR\#7},\; \id{at\_risk} = \id{false}\\[3pt]
\text{B: source of \id{refund} edited, no re-proof} & \id{done},\; \id{evidence\_ref} = \id{VR\#7},\; \id{at\_risk} = \id{true}\\
 & \quad\text{\footnotesize\itshape (``re-verify: source drifted''; still met, never silently demotes)}\\[3pt]
\text{C: superseding fresh proof appends} & \id{done},\; \id{evidence\_ref} = \id{VR\#12},\; \id{at\_risk} = \id{false}
\end{array}
\]
\caption{The non-monotonic \code{at\_risk} overlay: an obligation stays met when its witness goes stale after a source edit but is flagged for re-verification, and a superseding fresh proof clears the flag --- honesty without silent demotion.}
\label{fig:goal-at-risk}
\end{figure}

\subsection{Faithfulness --- the certified axis}\label{sec:goals-certified}

\code{met} and \code{governed} both sit \emph{downstream} of an unguarded seam: in the real workflow
the same agent that authors the definition of done then meets it, a textbook specification-gaming
setup. An agent can be perfectly honest at resolution --- never fabricate an artifact --- and still
``succeed'' by choosing a weak or incomplete definition of done (the ``wrong spec / vacuous proof''
problem). Faithfulness cannot be decided mechanically --- whether a formal acceptance faithfully
captures an informal intent is the formalization gap, and intent lives in a human's head --- so the
certified axis does not \emph{verify} faithfulness; it makes the seam visible, reviewed by a different
principal, coverage-checked, and hard to retrofit. Two mechanisms carry it. \textbf{Coverage:} each
obligation's \code{covers} field names the intent clauses it addresses, so an \code{intent\_clauses}
entry with no covering obligation surfaces as an uncovered-clause residual --- completeness becomes a
checkable projection rather than a judgement call. \textbf{Criteria review:} a \code{criteria\_review}
record signs off the acceptance \emph{set} --- ``this is what done means'' --- and its author
\emph{must} be a party other than the doer (\code{Human} or \code{Reviewer}, never the working
\code{Agent}). The certified verdict --- \code{faithfulness\_of} the goal --- is \code{Approved} when a
non-doer has signed off the definition \emph{and} no intent clause is left uncovered. A criterion that
\emph{bites} is evidence of faithfulness: a property refuted with a counterexample is stronger evidence
that the definition captures something real than an all-green suite of trivial edits. Certification also
sees the governance overrides: a goal certified with a formal-methods pack \emph{disabled} is a visibly
weaker claim than one certified with it enforced.

\subsection{Authorable today versus open}\label{sec:goals-status}

We are precise about what ships. The \emph{leaf atoms} of the property language already exist --- they
are the resolver helpers that compute \code{proved}, \code{has}, \code{no\_open\_gap}, freshness, and
contestation --- and the shipping authorable surface is the fixed \code{component}\,+\,\code{evidence}
struct, which is exactly the \code{exists}-over-one-component fragment those helpers evaluate. The
lineage-precise resolution rule of \S\ref{sec:goals-resolution}, the \code{at\_risk} overlay, and the
legacy-desugaring path are all shipped (in the reference resolver, both the Python and the JavaScript
implementations, held in parity by a harness). What is \emph{not yet built} is the full property-language
front-end that makes the language \emph{authorable} in its generality:

\begin{itemize}
\item a \emph{formula AST plus recursive evaluator} in place of the resolver's fixed \code{kind}-switch,
so \code{and}/\code{or}/\code{not}/\code{=>} and the role tags compose;
\item a \emph{component index} so \code{forall $f$ in glob("payments/**")}, \code{module($p$)}, and
\code{tag($l$)} can \emph{enumerate} their sets --- the one genuinely new piece of machinery, though it
indexes data (per-file source symbols, the high-stakes glob set) the trace already carries;
\item \emph{stable ids for \code{properties(S)}}, so a spec's authored verification goals are
addressable for \code{forall $p$ in properties(S)}.
\end{itemize}

\noindent Until then, goals are authored in the fixed struct (which desugars faithfully), and quantified
coverage is expressed through \code{[governed]} policies rather than \code{[met]} selectors. The design
of record is the two-sorted property language of \S\ref{sec:goals-criterion}; the evaluator for its
leaves is real; the combinator and enumeration layers are the honest outstanding work.
